\section{Trabalhos relacionados}

Esta seção apresenta uma revisão da literatura fundamentada em dois eixos
centrais para o posicionamento deste trabalho: a adoção de \textit{pipelines}
de ingestão adaptativos baseados em metadados e as soluções existentes para a
gestão e armazenamento de documentos fiscais. A partir dessa análise,
identificam-se as lacunas que justificam a
arquitetura proposta nesta pesquisa.

\subsection{\textit{Pipelines} de Ingestão de Dados Adaptativos}
\label{sec:pipelines_adaptativos}

A capacidade de se adaptar à evolução de esquemas é um dos principais
motivadores para a adoção de arquiteturas de ingestão de dados orientadas por
metadados. Chirumamilla (2024) aborda esse problema, ao
demonstrar como esquemas abstratos centralizados em repositórios de metadados,
permitem modificar modelos de dados, sem reescrever o código do
\textit{pipeline} de integração. Enquanto Chirumamilla (2024) aplica essa
elasticidade estrutural às mudanças em catálogos de produtos de comércio
eletrônico e varejo, nossa pesquisa direciona essa mesma premissa para o
cenário da escrituração fiscal.

No Brasil, os arquivos da EFD ICMS IPI possuem leiautes em arquivos texto
estruturados em árvore que sofrem revisões governamentais periódicas
\cite{nota_tecnica_efd}. O \textit{pipeline} de ingestão aqui proposto inova
ao utilizar esquemas JSON como contrato de metadados para interpretar, validar e
extrair os registros da EFD ICMS IPI dinamicamente. Dessa forma, quando a
legislação fiscal altera a estrutura de um bloco ou adiciona novos campos, a
adaptação do sistema requer apenas a atualização do artefato de mapeamento JSON,
conferindo alta manutenibilidade ao processamento de dados tributários.

\subsection{Armazenamento de documentos fiscais}
\label{sec:armazenamento_docs_fiscais}

No contexto da automação fiscal e da modernização dos sistemas de suporte ao
ecossistema do SPED, existem pesquisas que focam na descentralização e no ganho
de resiliência arquitetural para lidar com o grande volume de obrigações
tributárias. Nessa vertente, Tavares (2022) propõem uma
arquitetura baseada em microsserviços e contêineres \textit{Docker} para a
centralização, baixa e armazenamento de Notas Fiscais Eletrônicas (NF-e),
advindas da Secretaria da Fazenda (SEFAZ). O trabalho foca em solucionar os
gargalos operacionais da conferência manual e assegurar a guarda dos arquivos
XML (Linguagem de Marcação Extendida, tradução livre do inglês
\textit{Extensible Markup Language}) em um repositorio lógico de armazenamento
na nuvem, isolando o consumo e os recursos computacionais de cada empresa, para
prover alta disponibilidade e escalabilidade.

Embora Tavares (2022) demonstre a viabilidade de arquiteturas escaláveis para a
captura e preservação de documentos fiscais em nível de arquivo físico, a
solução limita-se à camada de aquisição e guarda em nuvem de arquivos já bem
estruturados em XML. A literatura ainda carece de abordagens robustas voltadas
para a ingestão, interpretação e estruturação granular de obrigações acessórias
textuais não tão bem estruturadas, como o EFD ICMS IPI. 

Nesse sentido, o presente trabalho avança e complementa a linha de investigação
de modernização contábil-tributária, ao propor um \textit{pipeline} de ingestão
dinâmico e orientado a metadados. Enquanto soluções como a de
\cite{tavares2022nfe} resolvem o problema do tráfego e do armazenamento bruto
do documento, a arquitetura aqui proposta foca no processamento relacional de
alta precisão, para assegurar a integridade dos dados em Bancos de Dados
Relacionais \textit{PostgreSQL}.